\section{Related Work}
\seclabel{related}
Within the Rocq ecosystem, an abundance of tactics enable algebraic simplification.
%
\citepos{boutin:reflection} ring
and field tactics\footnote{See the Rocq documentation:
  \url{https://rocq-prover.org/distrib/current/refman/addendum/ring.html} .}
let programmers discharge proof obligations involving (and requiring!)
addition, multiplication, and division operations.
\citepos{strub:coqmt} CoqMT extends Rocq's Calculus of
Inductive Constructions, allowing users to extend the conversion rule with
arbitrary decision procedures for first order theories (e.g. Presburger
arithmetic).
%
To ensure preservation of good meta-theoretical
properties, Strub extends only term level conversion.
%
Implementations of Hilbert's Nullstellensatz theorem
(\citepos{harrison:grobner} in HOL Light
and \citepos{pottier:grobner} in Rocq)
help users discharge proof obligations involving
polynomial equalities on a commutative integral domain.
%

Rocq's \SetoidRewrite{} is an advanced tactic library for setoid rewriting.%
\footnote{See the Rocq documentation:
  \url{https://rocq-prover.org/refman/addendum/generalized-rewriting.html}
  .} Disregarding the difference between the direct manipulation of
proof-terms in \idris{} and the tactic-based manipulation in Rocq,
\SetoidRewrite{} provides abstractions for manipulating parameterised
relations (covariant and contravariant), and users can register
setoids of interest and custom `morphisms' --- horn-like equational
clauses --- with the library. The various tactics in the library apply
these user-defined axioms to the goal. Users may also register
tactics, and the library includes an expressive collection of
term-traversal primitives (climbing up and down the syntax tree,
repeating sub-tactics, and so on). While \SetoidRewrite{} does not deal
with algebraic simplification directly, it may help in generalising
equality-based simplifiers to setoid-based simplifiers. In comparison,
\Frexlib{}'s setoid reasoning is minimal, implementing only the
necessary features for the library.

In Idris1,
\citeauthor{slama-brady:automatically-proving-equivalence-by-type-safe-reflection}~\citeyearpar{slama-brady:automatically-proving-equivalence-by-type-safe-reflection} and
\citeauthor{slama:thesis}~\citeyearpar{slama:thesis} implement a hierarchy of rewriting procedures for
algebraic structures of increasing complexity.  Though the procedures'
completeness is not enforced by type as in \Frexlib{}, these
simplifiers are based on a Knuth-Bendix resolution of critical pairs,
and so are likely to be complete.  \Frexlib{} also investigates a hierarchy of
rewriting procedures, but: (1)~frexlets are complete by
construction, (2) \Frexlib{} is based on normalisation-by-evaluation
(like Boutin's tactic, and unlike Slama and Brady's), and (3) \Frexlib{} is
extensible, with support for sufficiently motivated users to extend the library
with bespoke solvers.

Normalisation-by-evaluation is an established technique for
simplifying terms in a concrete equational theory, often involving
function types. One compelling example is \citepos{allais:nu-rules}
work, which demonstrates by a careful model
construction that the equational theory decided by
nor\-ma\-li\-sa\-tion-by-eval\-u\-a\-tion can be enriched with additional
rules. They implement a simply typed language internalising the
functorial and fusion laws
for list {\tt fold}, {\tt map}, and {\tt append} and prove
their construction sound and complete with respect to the extended
equational theory.

In Agda, \citepos{cockx:rewriting} and \citepos{cockx-et-al:taming-of-the-rew} `\texttt{-{}-rewriting}'
flag allows users to enrich the existing
reduction relation with new rules. Their implementation goes beyond Allais et al.'s:
it may restart stuck computations.
%
They leave to future work the soundness of user-provided reduction
rules, i.e.~ensuring rules neither introduce nontermination nor break
canonicity.
%
Unlike our approach, neither Allais et al's nor
Cockx et al's technique currently supports commutativity.

Existing formalisations for specific algebraic structures (monoids,
groups, rings, actions, etc.)
abound~\cite{mahboubi-tassi:mathcomp,agda-stdlib,hu-carette:agda-ct,lean-mathlib}. \Frexlib{}
is set apart from these by formalising the fragments of universal
algebra needed for its architecture. However, in \Frexlib{} we do not aspire
to formalise universal
algebra for its own sake.
Formalising more complete
fragments of the theory is an active area of research, with recent
contributions by \citet{GUNTHER2018147} and
\citet{DBLP:journals/corr/abs-2111-07936} in Agda.
\citet{carette-farmer-sharoda:leveraging-the-information-contained-in-theory-presentations}
generate \agda{} code for a comprehensive collection of multi-sorted
algebraic theories and their associated machinery via a pre-processing
phase from a much smaller description.
\citet{fiore-szamozvancev:soas-in-agda} similarly generate definitions
for the significantly more general second-order abstract syntax in
\agda{} via a preprocessing step, while formalising more of the
meta-theory as library code.  We consider it an open problem in this
domain to include the concise information as first-class data in the
meta-language while nonetheless enjoy the full ergonomics of
hand-written, inlined, definitions.

\agda{}'s category theory library~\cite{hu-carette:agda-ct} uses
setoids extensively. Each category includes both its homsets and their
dualisation. This choice allows for the operation `take the opposite
category' to be judgementally involutive. In turn, one can define
dualisation of a property purely formally, e.g. coproducts can be
defined as products in the opposite category. Unfortunately, this
design makes essential use of $\eta$-equality on records, yet
unsupported by \idris{}.

The Meta-F$\star$ language~\cite{martinez-et-al:metafstar} provides
normalisation tactics for commutative monoids and semi-rings through
its metaprogramming facilities. \Frexlib{}'s usage resembles
these tactics' usage, and we hope a
\Frexlib{} port to F$\star$ will make use of F$\star$'s metaprogramming
facilities to reduce some syntactic noise during goal extraction.

\section{Conclusions and Further Work}
\seclabel{conclusion}

We have presented a novel, mathematically structured, design for
algebraic simplification suites that guarantees sound and complete
simplification, even of user-defined simplifiers. Preliminary
evaluation shows that, despite a high level of abstraction, the
resulting library is responsive, with comparable functionality
to other libraries, in a combination of features that no single existing
library provides. \Frexlib{}'s unique design---the frex and
the fral---offer new prospects and questions.

Several computational type theories
support extensional function types and quotients, such as
observational type theory~\cite{altenkirch-et-al:ott-now,pujet-tabareau:ott-now-for-good,pujet-tabareau:impredicative-ott,pujet-leray-tabareau:ott-meets-cic} and cubical type theory~\cite{vezzosi-et-al:cubical-agda,cohen-et-al:cubical-tt}.
Extensional function types and quotient types are also typical reasons
for using setoids.  So it is natural to consider the tradeoffs for
implementing a similar library in a type theory with extensional
function types and quotient types.

Even in such rich type theories, we would still want to implement
setoid algebras.  Setoid algebras let us implement intensional aspects
such as proof printing, certification, and simplification, which we
view as useful functionality of a simplification library. Quotient
types cannot replace this functionality: eliminators of the quotient
must respect the quotient.

However, a richer type theory may make other aspects of the library
easier. For example, the commutative monoid fral and frex use the
power of an algebra by a setoid. If we worked in a type theory with
function extensionality, we could use the power of an algebra by a set
and do away with some of the setoid machinery involved.  A more
speculative direction involves using the improvements in the type
theory to implement category theoretic abstractions (presheaf
categories, ends and coends, 2-functors) which may allow us to `teach'
Idris not just universal algebra, but categorical algebra. We expand
on this idea further below.

Cubical type theories offer another potential implementation strategy
through non-trivial identity types.  Our Agda implementation
(\Fragmentlib) is compatible with Cubical Agda, as it type-checks with
the pragmas \texttt{-{}-safe} and \texttt{-{}-without-K}. Cubical identity types
turn types
into untruncated setoids, so our machinery applies unchanged.
However, cubical identity types cannot replace setoids.
Extracting intensional information is desirable and relevant  for cubical type theory too, so
we will still
implement setoid algebras. So cubical type theory will not
improve \Frexlib{}'s core design.

Implementing some of the internal data-structures may be
simplified by a more powerful type theory. For example, we can
implement the power of a model by an $n$-element set either as a
function or as an $n$-vector. It might be possible to use univalence
to ease transporting structure between these two different
implementations. However, using cubical abstractions---rather than
merely being compatible with a cubical type-theory---is a lot more
sophisticated. Such transport might lead to unexpected challenges, and
would require further experimentation.

\citepos{yallop-et-al:partially-static-data-as-frex} partial
evaluators include additional frexlets (e.g.~abelian groups,
semirings, distributive lattices). We plan to follow suit and port the
remaining simplifiers, then conduct larger evaluation and comparison studies.
The main challenge is that, unlike
\citeauthor{yallop-et-al:partially-static-data-as-frex}, we must
mechanising proofs showing these frexlets complete, which is costly.
One elegant motivation for including more
simplifiers is the following. The frex generalises the `ring of
polynomials over a ring' to that of an algebra of polynomials over an
algebra. By porting
\citeauthor{yallop-et-al:partially-static-data-as-frex}'s family of
representations, we will fully realise this generalisation.

Our experiment with reflection-based goal extraction as well as the reflection-based
interfaces of existing solvers show that with enough engineering
efforts, library designers can extract the goal equation from the goal
type. However, software engineers for dependently typed
languages are a scarce resource. We plan to explore other principled
approaches. In practice, when invoked inside a chain of equational
steps, the goal equation already appears in the source-code, albeit in
a context. Programmers seem willing to type the goal
equation once, since it documents the reasoning steps, but seem
unhappy to do so \emph{twice}. Generic
programming with holes\footnote{See Brad Hardy's Agda-Holes library:
  \url{https://github.com/bch29/agda-holes} .} may help.

Another promising direction is \emph{bootstrapping} of
the \Frexlib{} library using simplifier certification.
%
Bootstrapping might start with a hierarchy of inefficient simplifiers
that are easy to implement.
%
Next, these simplifiers may then be used to develop a hierarchy of
more efficient simplifiers and proof-simplifiers.
%
Finally, the certification mechanism can extract proofs to complete
the bootstrap.

We would also like to extend \Frexlib{}'s design beyond algebraic
structures. More general notions of theories abound: multisorted,
second-order/parameterised, and essentially algebraic.
%
Supporting these may allow \Frexlib{} to cover much more complex situations, such as decision
procedures for first order theories (e.g.~Presburger arithmetic,
cf.~\citepos{strub:coqmt} CoqMT) normalisation-by-evaluation
for fusion laws~\cite{allais:nu-rules}, and equational manipulation of
big-operators~\cite{betot-et-al:big-operators,markert:big-operators,lau:big-operators}.
%
Note that \Frexlib{} can already deal with big-operators such as \IdrisFunction{sum} so
long as the argument list is a concrete collection of constants and
variables such as
\IdrisFunction{sum}\ \IdrisData{[2,}\ \IdrisBound{x}\IdrisData{]}. We
only need the more sophisticated theories when the length of the lists
is abstract.

\Frexlib{} uses many
category-theoretic concepts, but the library itself is oblivious to
category theory. %% Some care is needed for such a category theory library: it
%% needs to avoid postulates that block reduction.
We hope that making use of a rich category theory library like \citepos{hu-carette:agda-ct} {\tt agda-categories} might
lead to a sleeker and even more modular \Frexlib{} implementation. More
specifically, we plan to explore a general treatment of involutive
algebras following \citeauthor{jacobs:involution}~\citeyearpar{jacobs:involution}, and Power's
\emph{distributive tensor} of equational
theories~\cite{power:discrete-lawvere-theories,%
  hyland-power:discrete-lawvere-theories}
for a uniform treatment of semi-ring varieties.
Instantiating each of the $6$ semi-group varieties makes it possible to cover
each instance of the following combinations:
\[
\left(
\begin{array}{@{}l@{}}\hphantom{\cup{}}
  \set{\text{ordinary}}\times\set{\text{ordinary}, \text{involutive},
        \text{non-reversing involutive}}
\\\cup
\set{\text{commutative}}\times\set{\text{ordinary}, \text{involutive}
}
\end{array}
\right)\times\set{\begin{array}{@{}c@{}}
    \text{commutative}\\
    \text{semigroup}, \text{monoid}, \text{group}
  \end{array}
}
\]
and modularly construct $(2+3)\times 3 = 15$ semi-ring varieties,
including rings and semirings. As this example shows, this kind of
modular treatment can provide a multiplicative development boost.

\newpage

%%

%% Also, units of
%% measurement.




%% NbE. Proof-simplification.


%% Morita equivalence, cat theory
%% Fusion laws


%% Things that will require significant innovation to fit into frex:

%% Presburger arithmetic and Nullstellensatz
%% Arbitrary rewriting
%% Incomplete decision procedure

%% Other languages

%% Other frexy directions:

%% Indexing-modulo-equations

%% Holey reasoning with zippers
%% (instead of elaborator refelction/macros, use overloading to construct the LHS/RHS + context)

%% Future of staging in Idris, if we want the reflection magic.

%% WG1: interoperability

%% Exciting initiatives for cross-prover between provers/dependently
%% typed systems, eg Euro-Proof-Network

%% core-dump:
%% %
